\documentclass[12pt]{article}

\usepackage[a4paper, total={6in, 8in}]{geometry}
\usepackage{textcomp}

\begin{document}
\setlength{\parindent}{0pt}

\def \t {\textrightarrow}
\def \v {\vspace{1em}}
\def \bi {\begin{itemize}}
\def \ei {\end{itemize}}
\def \s[#1] {\section*{#1}}
\def \ss[#1] {\subsection*{#1}}
\def \sss[#1] {\subsubsection*{#1}}

\s[Bergson]
Lui fonda la liberta sulla durata \t azioni libere sono libere solo se sono coerenti con la sua natura \\
Noi compiamo azioni abitudinarie, che sono quotidiane che non sono negative pero certamente non libere \\
Se esiste liberta o libero arbitrio \t quesito è malposto \t noi dobbiamo considerare la coscienza non in termini meccanici o deterministici \t o/o non appartiene alla coscienza, categorizzaizone troppo rigida che non appartiene alla coscienza \\
Det. vs lib. arb. \t nella coscienza è piu complicato, gomitolo, quindi quesisto è posto male

\ss[Materia e memoria]
Il tempo spazializzato si contrappone alla durata \t e questo si ripercuote in un altra contrapposizione tra realta meccanica esterna e meccanica interna, che è massimamente creativa \\
Qua ripronpone dualismo tra realta meccanica e realta della coscienza \t e pone la questione del passaggio da un all altra \\
Questo passaggio è possibile perche ha evidenziato come a volte la coscienza possa solidificarsi in azioni ripetitive, meccaniche \t ha visto che il meccanismo della realta fuori puo influenzarmi a tal punto da rende la coscienza ripetitiva, meccanica \\
C'è possibilita di influenza \t quindi bisogna vedere se il passaggio avviene, come avviene e se ha ripercussioni \\
Trova collegamento tra materia (r. esterna) e memoria (realta interna) \\
Fa considerazioni simili a quelle di Freud \t c'è tangenza per alcuni aspetti, anche se versanti diversi (no psicanalisi qua) \\
Distingue coscienza e cervello \t nella coscienza c'è infinitamente di piu del cervello \\
Distingue 3 momenti della coscienza

\sss[Memoria]
Coincide con tutta la storia della coscienza \t il passato che mi porto dietro di cui pero non sono consapevole \t io sono qui perche ho vissuto una serie di fatti ed eventi che mi hanno portato ad essere chi sono \\
Il passato sta nella persona, anche se in modo inconcio (non ci ricordiamo ogni dettaglio del nostro vissuto) \\
La memoria è il gomitolo

\sss[Ricordo]
È quello che viene ripescato dalla memoria ed è cio che mi serve nel presente \t è selezione di qualcosa che c'è nella memoria, che viene scelto perche serve per orientarsi nel presente \\
Ricordo è una minima parte del mio passato \\
L'opera di selezione viene fatta dal cervello \t è una azione corporea \t seleziona secondo un criterio di utilita \\
Nel cervello passa quindi una piccola parte di cio che c'è nella coscienza \t nel cervello ci sono ricordi, ma è una selezione di cio che è la memoria (che è la coscienza)

\sss[Percezione]
La percezione permette di muovermi all interno del mondo \t permette al corpo di essere tra altri corpi, e di muovermi nella realta sulla base delle mie esperienze passate \t i ricordi orientano la percezione \\
La mia memoria, per realizzarsi, ha bisogno del corpo, perche la percezione è fisica \t quindi pero la mia memoria è indipendente dal mio corporea

\v

Memoria = coscienza, percezione = corpo \t la percezione nasce sulla base di una selezione = il mio modo di stare nel mondo è orientato non dal modo in cui penso di pormi oggi, ma dal passato che mi ha costruito (con cui devo fare dei conti) \\
La liberta ha a che fare con questo discorso \t io sto nel mondo con la percezione, e la mia liberta viene limitata dalla percezione \t io sono libero a stare nel mondo a partire da cio che sono \t cio che sono stato e cio che mi porto sta influenzando il modo in cui mi pongo nella realta \t non sono massimamente libero quindi \\
Collegamento con Freud è l'inconscio \t ma analisi clinica li 

\v

Non esiste quindi dualismo tra res cogitans ed extensa \t no Cartesio \t spirito e memoria non sono due cose lontane, sono due aspetti della stessa cosa

\ss[Evoluzione creatrice]
Evoluzionismo cosmologico \\
Parla facendo un analisi e critica delle teorie evoluzionistiche finora proposte, che si dividono in:
\bi
  \item teorie meccanicistiche = Darwin \t evoluzione che spiega la realta in termine di causalita efficiente
  \item teorie finalistiche \t ev. si determina sulla base del futuro \t mi evolvo perche devo realizzare uno scopo
\ei
Queste sono entrambe teorie deterministiche \t perche ripropongono schema: o ev. nel passato o nel futuro ? \\
La vita pero cosi non si evolve \t la vita biologica non è una macchina, non si ripete \t l evoluzione della vita è incessante novita, impredevibilita \\
L'evoluzione di aggroviglia su se stessa \t ed è caratterizzata dallo slancio vitale, che si è fermato con la nascita della materia \\
L'evoluzione è come una bomba che esplode, i cui frammenti vanno in tutte le direzioni \t la vita biologica è imprevedibile \\
La traiettoria evolutiva non è lineare \\
L'evoluzione è andata in modo divergente \t i viventi si sono specializzati in modi molto distinti (es. piante vs animali) \t anche animali sono andati verso l'istinto o intelligenza \\
Nella materia lo slancio vitale si è degradato \t l'unico ente materiale in cui è ancora presente ma in modo diverso è l uomo, che ha mantenuto impredevibilita (Hume, componente legata all'intelligenza) e spirito creatrice

\v

Slancio è metafisico? \\
Nell'uomo pero c'è anche istinto \t animale solo istinto \\
L'istinto è la facolta di usare strumenti organici e di costruirli (es. orso che usa l'albero o gli uccelli che costruiscono nido) \\
L'istinto funziona, offre soluzioni adeguate (per un solo problema pero: risponde solo a quella domanda, non puo variare perche è ripetitivo) \\
L'istinto puo conoscere le cose ma non i rapporti tra le cose \\
L'intelligenza è invece la facolta di fabbricare, utilizzare strumenti inorganici \t puo crare strumenti artificiali che non esistono in natura \\
Inoltre intelligenza non è ereditaria, mentre istinto si \t intelligenza è creativa, trova piu soluzioni a un problema \\
Come è passato uomo da istinto a intelligenza? \\
Intelligenza stabilisce nessi tra enti ed è cosciente \\
Ne l'istinto da solo ne intelligenza da sola fanno pero entrare nella realta \t ci fanno entrare nella realta l'intuizione = istinto + intelligenza \\
Int. non è mai totalmente separata dall'istinto \t uomo è anche animale \t l'ingelligenza fa girare intorno all'oggetto, ma poi manca l'istinto che permette intuizione \\
Non si entra solo in ottica logico-razionale con l'oggeto, ma entra aspetto intuitivo?? \\
L'intuizione è un processo reale? \t si, ed è per eccellenza nell'intuizione estetica \t colgo l'oggetto come completamente sviconlato dalla realta (es. l opera d'arte) \\
Dice anche che l'intuizione è l'organo della metafisica \t per entrare in contatto con dimensione non sensibile serve intuizione \t entro in contato con essenza degli enti e con la durata, che non è sensibile
\end{document}
